% !TEX root = ../main.tex
\chapter{Nettoyage et Pr\'eparation des Donn\'ees}
\label{ch:nettoyage}

\begin{intuition}
On estime que les data scientists consacrent environ 80\,\% de leur temps au nettoyage et \`a la pr\'eparation des donn\'ees. Des donn\'ees mal pr\'epar\'ees m\`enent in\'evitablement \`a des mod\`eles erron\'es : <<~Garbage in, garbage out. ~>>
\end{intuition}

% ============================================================
\section{Pourquoi nettoyer les donn\'ees ?}
\index{nettoyage des donn\'ees}

Les donn\'ees r\'eelles sont rarement propres. On rencontre fr\'equemment :
\begin{itemize}
    \item des \textbf{valeurs manquantes}\index{valeurs manquantes} (capteurs d\'efaillants, champs non remplis) ;
    \item des \textbf{doublons}\index{doublons} (saisies multiples, fusions de tables) ;
    \item des \textbf{types incorrects}\index{types de donn\'ees} (nombres stock\'es comme texte) ;
    \item des \textbf{valeurs aberrantes}\index{valeurs aberrantes} (erreurs de saisie) ;
    \item des \textbf{incoh\'erences textuelles} (majuscules, espaces, abr\'eviations).
\end{itemize}

% ============================================================
\section{Pipeline de nettoyage}

\begin{center}
\begin{tikzpicture}[
    node distance=0.9cm,
    block/.style={rectangle, draw=teal!80, fill=teal!10, rounded corners, text width=4cm, align=center, minimum height=0.9cm, font=\small},
    arrow/.style={->, thick, >=stealth, color=gray!70}
]
\node[block] (raw) {Donn\'ees brutes};
\node[block, below=of raw] (dup) {Suppression des doublons};
\node[block, below=of dup] (types) {Correction des types};
\node[block, below=of types] (miss) {Traitement des manquants};
\node[block, below=of miss] (out) {Traitement des outliers};
\node[block, below=of out] (text) {Nettoyage textuel};
\node[block, below=of text] (feat) {Feature engineering};
\node[block, below=of feat] (enc) {Encodage et mise \`a l'\'echelle};
\node[block, below=of enc] (clean) {Donn\'ees propres};

\foreach \a/\b in {raw/dup, dup/types, types/miss, miss/out, out/text, text/feat, feat/enc, enc/clean}
    \draw[arrow] (\a) -- (\b);
\end{tikzpicture}
\end{center}

% ============================================================
\section{\'Etude de cas : le jeu Titanic}
\index{Titanic}

\begin{codeblock}[title=\textbf{Python}]
\begin{minted}{python}
import pandas as pd
import numpy as np
import seaborn as sns

titanic = sns.load_dataset('titanic')
print(f"Dimensions initiales : {titanic.shape}")
print(titanic.info())
\end{minted}
\end{codeblock}

\begin{codeoutput}
\begin{minted}{text}
Dimensions initiales : (891, 15)
<class 'pandas.core.frame.DataFrame'>
RangeIndex: 891 entries, 0 to 890
Data columns (total 15 columns):
 #   Column       Non-Null Count  Dtype
 0   survived     891 non-null    int64
 1   pclass       891 non-null    int64
 2   sex          891 non-null    object
 3   age          714 non-null    float64
 ...
 7   embarked     889 non-null    object
 11  deck         203 non-null    category
\end{minted}
\end{codeoutput}

% ============================================================
\section{Valeurs manquantes}
\index{valeurs manquantes}\index{MCAR}\index{MAR}\index{MNAR}

\subsection{Types de donn\'ees manquantes}

\begin{definition}[Classification des valeurs manquantes]
\begin{itemize}
    \item \textbf{MCAR} (\emph{Missing Completely At Random}) : le caract\`ere manquant ne d\'epend d'aucune variable.
    \item \textbf{MAR} (\emph{Missing At Random}) : le caract\`ere manquant d\'epend d'autres variables observ\'ees.
    \item \textbf{MNAR} (\emph{Missing Not At Random}) : le caract\`ere manquant d\'epend de la valeur elle-m\^eme.
\end{itemize}
\end{definition}

\begin{exemple}
Dans le Titanic, les cabines (\py{deck}) manquent surtout pour les passagers de 3\`eme classe (MAR). L'\^age pourrait \^etre MCAR ou MAR selon que les passagers sans billet complet n'avaient pas d'\^age enregistr\'e.
\end{exemple}

\subsection{D\'etection et traitement}

\begin{codeblock}[title=\textbf{Python}]
\begin{minted}{python}
# Détection
print(titanic.isnull().sum().sort_values(ascending=False).head(5))

# Stratégie 1 : supprimer la colonne deck (trop de manquants)
titanic_clean = titanic.drop(columns=['deck'])

# Stratégie 2 : remplir age par la médiane par classe
titanic_clean['age'] = titanic_clean.groupby('pclass')['age'] \
    .transform(lambda x: x.fillna(x.median()))

# Stratégie 3 : remplir embarked par le mode
titanic_clean['embarked'].fillna(
    titanic_clean['embarked'].mode()[0], inplace=True)
titanic_clean['embark_town'].fillna(
    titanic_clean['embark_town'].mode()[0], inplace=True)

print(f"\nManquants restants : {titanic_clean.isnull().sum().sum()}")
\end{minted}
\end{codeblock}

\begin{codeoutput}
\begin{minted}{text}
deck           688
age            177
embarked         2
embark_town      2
survived         0

Manquants restants : 0
\end{minted}
\end{codeoutput}

\begin{bonnepratique}
L'imputation par la m\'ediane \emph{conditionn\'ee} (ici par classe) est pr\'ef\'erable \`a la m\'ediane globale car elle pr\'eserve la structure des sous-groupes. Pour des donn\'ees temporelles, utilisez \py{ffill()} ou \py{bfill()}.
\end{bonnepratique}

% ============================================================
\section{Doublons}
\index{doublons}

\begin{codeblock}[title=\textbf{Python}]
\begin{minted}{python}
# Détection
n_dup = titanic_clean.duplicated().sum()
print(f"Nombre de doublons : {n_dup}")

# Inspection des doublons éventuels
if n_dup > 0:
    print(titanic_clean[titanic_clean.duplicated(keep=False)]
          .sort_values(by='name').head())
    titanic_clean = titanic_clean.drop_duplicates()
    print(f"Après suppression : {titanic_clean.shape}")
\end{minted}
\end{codeblock}

\begin{codeoutput}
\begin{minted}{text}
Nombre de doublons : 0
\end{minted}
\end{codeoutput}

% ============================================================
\section{Conversion de types}
\index{conversion de types}

\begin{codeblock}[title=\textbf{Python}]
\begin{minted}{python}
# Convertir survived en booléen
titanic_clean['survived'] = titanic_clean['survived'].astype(bool)

# Convertir pclass en catégoriel ordonné
titanic_clean['pclass'] = pd.Categorical(
    titanic_clean['pclass'], categories=[1, 2, 3], ordered=True)

# Exemple de conversion de date (sur données fictives)
dates_str = pd.Series(['2023-01-15', '2023-02-20', '2023-03-10'])
dates = pd.to_datetime(dates_str)
print(dates.dt.month)
\end{minted}
\end{codeblock}

\begin{codeoutput}
\begin{minted}{text}
0    1
1    2
2    3
dtype: int32
\end{minted}
\end{codeoutput}

% ============================================================
\section{Traitement des valeurs aberrantes}
\index{valeurs aberrantes}

\begin{codeblock}[title=\textbf{Python}]
\begin{minted}{python}
# Méthode IQR pour le tarif
Q1 = titanic_clean['fare'].quantile(0.25)
Q3 = titanic_clean['fare'].quantile(0.75)
IQR = Q3 - Q1

borne_inf = Q1 - 1.5 * IQR
borne_sup = Q3 + 1.5 * IQR

# Option 1 : écrêtage (capping/winsorizing)
titanic_clean['fare_capped'] = titanic_clean['fare'].clip(
    lower=borne_inf, upper=borne_sup)

print(f"Avant : max = {titanic_clean['fare'].max():.2f}")
print(f"Après écrêtage : max = {titanic_clean['fare_capped'].max():.2f}")
print(f"Borne supérieure : {borne_sup:.2f}")
\end{minted}
\end{codeblock}

\begin{codeoutput}
\begin{minted}{text}
Avant : max = 512.33
Après écrêtage : max = 65.63
Borne supérieure : 65.63
\end{minted}
\end{codeoutput}

\begin{attention}
Ne supprimez jamais les outliers sans r\'eflexion. Un tarif de 512\,\$ peut \^etre l\'egitime (suite de luxe). L'\'ecr\^etage est souvent pr\'ef\'erable \`a la suppression. Documentez toujours vos choix.
\end{attention}

% ============================================================
\section{Nettoyage textuel}
\index{nettoyage textuel}

\begin{codeblock}[title=\textbf{Python}]
\begin{minted}{python}
# Exemple avec des données textuelles
noms = pd.Series(['  Alice ', 'BOB', 'charlie', ' Dave  '])
print("Avant :", noms.tolist())

noms_clean = (noms.str.strip()
                  .str.lower()
                  .str.capitalize())
print("Après :", noms_clean.tolist())

# Remplacement de motifs
villes = pd.Series(['New-York', 'new york', 'NEW YORK', 'N.Y.'])
villes_norm = villes.str.lower().str.replace(
    r'[^a-z\s]', '', regex=True).str.strip()
print("Villes normalisées :", villes_norm.tolist())
\end{minted}
\end{codeblock}

\begin{codeoutput}
\begin{minted}{text}
Avant : ['  Alice ', 'BOB', 'charlie', ' Dave  ']
Après : ['Alice', 'Bob', 'Charlie', 'Dave']
Villes normalisées : ['newyork', 'new york', 'new york', 'ny']
\end{minted}
\end{codeoutput}

% ============================================================
\section{Feature engineering}
\index{feature engineering}

\begin{codeblock}[title=\textbf{Python}]
\begin{minted}{python}
# Créer de nouvelles variables à partir des existantes
titanic_clean['family_size'] = (titanic_clean['sibsp']
                                + titanic_clean['parch'] + 1)
titanic_clean['is_child'] = titanic_clean['age'] < 18
titanic_clean['fare_per_person'] = (titanic_clean['fare']
                                    / titanic_clean['family_size'])

# Discrétiser l'âge en catégories
titanic_clean['age_group'] = pd.cut(
    titanic_clean['age'],
    bins=[0, 12, 18, 35, 60, 100],
    labels=['Enfant', 'Ado', 'Jeune adulte', 'Adulte', 'Senior'])

print(titanic_clean['age_group'].value_counts())
\end{minted}
\end{codeblock}

\begin{codeoutput}
\begin{minted}{text}
Jeune adulte    371
Adulte          267
Enfant           89
Ado              79
Senior           85
Name: age_group, dtype: int64
\end{minted}
\end{codeoutput}

% ============================================================
\section{Encodage des variables cat\'egorielles}
\index{encodage}\index{one-hot encoding}\index{label encoding}

\begin{codeblock}[title=\textbf{Python}]
\begin{minted}{python}
# One-hot encoding avec pandas
embarked_dummies = pd.get_dummies(
    titanic_clean['embarked'], prefix='embarked', dtype=int)
print(embarked_dummies.head())

# Label encoding avec sklearn
from sklearn.preprocessing import LabelEncoder

le = LabelEncoder()
titanic_clean['sex_encoded'] = le.fit_transform(
    titanic_clean['sex'])
print("\nMapping :", dict(zip(le.classes_,
                              le.transform(le.classes_))))
\end{minted}
\end{codeblock}

\begin{codeoutput}
\begin{minted}{text}
   embarked_C  embarked_Q  embarked_S
0           0           0           1
1           1           0           0
2           0           0           1
3           0           0           1
4           0           0           1

Mapping : {'female': 0, 'male': 1}
\end{minted}
\end{codeoutput}

\begin{attention}
Le \textbf{label encoding}\index{label encoding} impose un ordre artificiel aux cat\'egories. Il ne convient que pour les variables ordinales ou les algorithmes \`a base d'arbres. Pour les mod\`eles lin\'eaires, pr\'ef\'erez le \textbf{one-hot encoding}\index{one-hot encoding}.
\end{attention}

% ============================================================
\section{Mise \`a l'\'echelle}
\index{standardisation}\index{normalisation}\index{StandardScaler}\index{MinMaxScaler}

\begin{definition}[Standardisation et normalisation]
\begin{itemize}
    \item \textbf{Standardisation} (StandardScaler) : $z = \dfrac{x - \mu}{\sigma}$, r\'esultat de moyenne 0 et d'\'ecart-type 1.
    \item \textbf{Normalisation} (MinMaxScaler) : $x' = \dfrac{x - x_{\min}}{x_{\max} - x_{\min}}$, r\'esultat dans $[0, 1]$.
\end{itemize}
\end{definition}

\begin{codeblock}[title=\textbf{Python}]
\begin{minted}{python}
from sklearn.preprocessing import StandardScaler, MinMaxScaler

cols = ['age', 'fare']
data_num = titanic_clean[cols].copy()

# Standardisation
scaler_std = StandardScaler()
data_std = pd.DataFrame(scaler_std.fit_transform(data_num),
                        columns=[c + '_std' for c in cols])

# Normalisation
scaler_mm = MinMaxScaler()
data_norm = pd.DataFrame(scaler_mm.fit_transform(data_num),
                         columns=[c + '_norm' for c in cols])

print("Standardisé :")
print(data_std.describe().loc[['mean', 'std']].round(2))
print("\nNormalisé :")
print(data_norm.describe().loc[['min', 'max']].round(2))
\end{minted}
\end{codeblock}

\begin{codeoutput}
\begin{minted}{text}
Standardisé :
      age_std  fare_std
mean     0.00      0.00
std      1.00      1.00

Normalisé :
      age_norm  fare_norm
min       0.00       0.00
max       1.00       1.00
\end{minted}
\end{codeoutput}

\begin{bonnepratique}
Appliquez la mise \`a l'\'echelle \emph{apr\`es} la s\'eparation train/test pour \'eviter la fuite d'information (\emph{data leakage}\index{data leakage}). Utilisez \py{fit\_transform()} sur le train et \py{transform()} sur le test.
\end{bonnepratique}

% ============================================================
\section{Exercices}

\begin{exercice}[$\star$]
Chargez le jeu \py{titanic}. Affichez le nombre et le pourcentage de valeurs manquantes par colonne. Quelles colonnes ont plus de 50\,\% de manquants ?
\end{exercice}

\begin{exercice}[$\star$]
Supprimez les doublons du jeu \py{titanic} (s'il y en a). Puis convertissez la colonne \py{sex} en variable cat\'egorielle avec \py{pd.Categorical()}.
\end{exercice}

\begin{exercice}[$\star\star$]
Imputez les valeurs manquantes de \py{age} par la m\'ediane conditionn\'ee par \py{sex} et \py{pclass} (6 groupes). Comparez la distribution de l'\^age avant et apr\`es imputation avec des histogrammes superpos\'es.
\end{exercice}

\begin{exercice}[$\star\star$]
Cr\'eez les variables suivantes : \py{title} (extrait du nom avec une expression r\'eguli\`ere), \py{is\_alone} (vrai si \py{family\_size} = 1), \py{fare\_bin} (quartiles du tarif avec \py{pd.qcut}). Calculez le taux de survie pour chacune de ces nouvelles variables.
\end{exercice}

\begin{exercice}[$\star\star\star$]
Construisez un pipeline complet de nettoyage pour le jeu \py{titanic} : (a) suppression de \py{deck}, (b) imputation de \py{age} et \py{embarked}, (c) cr\'eation de \py{family\_size} et \py{is\_child}, (d) one-hot encoding de \py{sex}, \py{embarked}, \py{class}, (e) standardisation des variables num\'eriques. Le r\'esultat final doit \^etre un DataFrame enti\`erement num\'erique sans valeurs manquantes.
\end{exercice}

\begin{exercice}[$\star\star\star$]
Chargez le jeu \py{penguins} (\py{sns.load\_dataset('penguins')}). Ce jeu contient des valeurs manquantes. Appliquez trois strat\'egies diff\'erentes (suppression, imputation par m\'ediane, imputation par m\'ediane conditionn\'ee par esp\`ece). Comparez les statistiques descriptives obtenues avec chaque strat\'egie et discutez de la meilleure approche.
\end{exercice}

% ============================================================
\begin{formules}
\textbf{R\'esum\'e du chapitre -- Nettoyage et pr\'eparation}
\begin{itemize}
    \item \textbf{Manquants} : \py{isnull().sum()}, types MCAR/MAR/MNAR, \py{fillna()}, \py{dropna()}.
    \item \textbf{Doublons} : \py{duplicated()}, \py{drop\_duplicates()}.
    \item \textbf{Types} : \py{astype()}, \py{pd.to\_datetime()}, \py{pd.Categorical()}.
    \item \textbf{Outliers} : m\'ethode IQR, \'ecr\^etage avec \py{clip()}.
    \item \textbf{Texte} : \py{str.strip()}, \py{str.lower()}, \py{str.replace()}.
    \item \textbf{Feature engineering} : \py{pd.cut()}, \py{pd.qcut()}, combinaison de colonnes.
    \item \textbf{Encodage} : \py{pd.get\_dummies()} (one-hot), \py{LabelEncoder} (ordinal).
    \item \textbf{Mise \`a l'\'echelle} : $z = \frac{x - \mu}{\sigma}$ (StandardScaler), $x' = \frac{x - x_{\min}}{x_{\max} - x_{\min}}$ (MinMaxScaler).
    \item \textbf{R\`egle d'or} : documenter chaque transformation et \'eviter le \emph{data leakage}.
\end{itemize}
\end{formules}
